% Options for packages loaded elsewhere
\PassOptionsToPackage{unicode}{hyperref}
\PassOptionsToPackage{hyphens}{url}
\PassOptionsToPackage{dvipsnames,svgnames,x11names}{xcolor}
%
\documentclass[
  11pt,
]{article}
\usepackage{amsmath,amssymb}
\usepackage{iftex}
\ifPDFTeX
  \usepackage[T1]{fontenc}
  \usepackage[utf8]{inputenc}
  \usepackage{textcomp} % provide euro and other symbols
\else % if luatex or xetex
  \usepackage{unicode-math} % this also loads fontspec
  \defaultfontfeatures{Scale=MatchLowercase}
  \defaultfontfeatures[\rmfamily]{Ligatures=TeX,Scale=1}
\fi
\usepackage{lmodern}
\ifPDFTeX\else
  % xetex/luatex font selection
\fi
% Use upquote if available, for straight quotes in verbatim environments
\IfFileExists{upquote.sty}{\usepackage{upquote}}{}
\IfFileExists{microtype.sty}{% use microtype if available
  \usepackage[]{microtype}
  \UseMicrotypeSet[protrusion]{basicmath} % disable protrusion for tt fonts
}{}
\makeatletter
\@ifundefined{KOMAClassName}{% if non-KOMA class
  \IfFileExists{parskip.sty}{%
    \usepackage{parskip}
  }{% else
    \setlength{\parindent}{0pt}
    \setlength{\parskip}{6pt plus 2pt minus 1pt}}
}{% if KOMA class
  \KOMAoptions{parskip=half}}
\makeatother
\usepackage{xcolor}
\usepackage[margin=2.4cm]{geometry}
\usepackage{longtable,booktabs,array}
\usepackage{calc} % for calculating minipage widths
% Correct order of tables after \paragraph or \subparagraph
\usepackage{etoolbox}
\makeatletter
\patchcmd\longtable{\par}{\if@noskipsec\mbox{}\fi\par}{}{}
\makeatother
% Allow footnotes in longtable head/foot
\IfFileExists{footnotehyper.sty}{\usepackage{footnotehyper}}{\usepackage{footnote}}
\makesavenoteenv{longtable}
\setlength{\emergencystretch}{3em} % prevent overfull lines
\providecommand{\tightlist}{%
  \setlength{\itemsep}{0pt}\setlength{\parskip}{0pt}}
\setcounter{secnumdepth}{-\maxdimen} % remove section numbering
\usepackage{newunicodechar}
\usepackage{amssymb}
\newunicodechar{±}{\ensuremath{\pm}}
\newunicodechar{²}{\ensuremath{^{2}}}
\newunicodechar{³}{\ensuremath{^{3}}}
\newunicodechar{·}{\ensuremath{\cdot}}
\newunicodechar{¹}{\ensuremath{^{1}}}
\newunicodechar{×}{\ensuremath{\times}}
\newunicodechar{ő}{\H{o}}
\newunicodechar{Σ}{\ensuremath{\Sigma}}
\newunicodechar{α}{\ensuremath{\alpha}}
\newunicodechar{β}{\ensuremath{\beta}}
\newunicodechar{δ}{\ensuremath{\delta}}
\newunicodechar{η}{\ensuremath{\eta}}
\newunicodechar{θ}{\ensuremath{\theta}}
\newunicodechar{ι}{\ensuremath{\iota}}
\newunicodechar{κ}{\ensuremath{\kappa}}
\newunicodechar{λ}{\ensuremath{\lambda}}
\newunicodechar{μ}{\ensuremath{\mu}}
\newunicodechar{ν}{\ensuremath{\nu}}
\newunicodechar{π}{\ensuremath{\pi}}
\newunicodechar{ρ}{\ensuremath{\rho}}
\newunicodechar{σ}{\ensuremath{\sigma}}
\newunicodechar{τ}{\ensuremath{\tau}}
\newunicodechar{ψ}{\ensuremath{\psi}}
\newunicodechar{ω}{\ensuremath{\omega}}
\newunicodechar{…}{\ldots{}}
\newunicodechar{′}{\ensuremath{'}}
\newunicodechar{″}{\ensuremath{''}}
\newunicodechar{⁰}{\ensuremath{^{0}}}
\newunicodechar{⁴}{\ensuremath{^{4}}}
\newunicodechar{⁵}{\ensuremath{^{5}}}
\newunicodechar{⁶}{\ensuremath{^{6}}}
\newunicodechar{⁷}{\ensuremath{^{7}}}
\newunicodechar{⁸}{\ensuremath{^{8}}}
\newunicodechar{⁹}{\ensuremath{^{9}}}
\newunicodechar{⁺}{\ensuremath{^{+}}}
\newunicodechar{⁻}{\ensuremath{^{-}}}
\newunicodechar{₀}{\ensuremath{_{0}}}
\newunicodechar{₁}{\ensuremath{_{1}}}
\newunicodechar{₂}{\ensuremath{_{2}}}
\newunicodechar{₃}{\ensuremath{_{3}}}
\newunicodechar{₄}{\ensuremath{_{4}}}
\newunicodechar{₅}{\ensuremath{_{5}}}
\newunicodechar{₆}{\ensuremath{_{6}}}
\newunicodechar{₇}{\ensuremath{_{7}}}
\newunicodechar{₈}{\ensuremath{_{8}}}
\newunicodechar{ℓ}{\ensuremath{\ell}}
\newunicodechar{ℛ}{\ensuremath{\mathcal{R}}}
\newunicodechar{→}{\ensuremath{\to}}
\newunicodechar{↦}{\ensuremath{\mapsto}}
\newunicodechar{∅}{\ensuremath{\emptyset}}
\newunicodechar{∈}{\ensuremath{\in}}
\newunicodechar{∉}{\ensuremath{\notin}}
\newunicodechar{∎}{\ensuremath{\blacksquare}}
\newunicodechar{∏}{\ensuremath{\prod}}
\newunicodechar{−}{\ensuremath{-}}
\newunicodechar{∖}{\ensuremath{\setminus}}
\newunicodechar{√}{\ensuremath{\surd}}
\newunicodechar{∞}{\ensuremath{\infty}}
\newunicodechar{∤}{\ensuremath{\nmid}}
\newunicodechar{∥}{\ensuremath{\parallel}}
\newunicodechar{∩}{\ensuremath{\cap}}
\newunicodechar{∪}{\ensuremath{\cup}}
\newunicodechar{≈}{\ensuremath{\approx}}
\newunicodechar{≠}{\ensuremath{\neq}}
\newunicodechar{≡}{\ensuremath{\equiv}}
\newunicodechar{≢}{\ensuremath{\not\equiv}}
\newunicodechar{≤}{\ensuremath{\le}}
\newunicodechar{≥}{\ensuremath{\ge}}
\newunicodechar{⊄}{\ensuremath{\not\subseteq}}
\newunicodechar{⊆}{\ensuremath{\subseteq}}
\newunicodechar{⊔}{\ensuremath{\sqcup}}
\newunicodechar{⌊}{\ensuremath{\lfloor}}
\newunicodechar{⌋}{\ensuremath{\rfloor}}
\newunicodechar{𝒞}{\ensuremath{\mathcal{C}}}
\newunicodechar{𝒫}{\ensuremath{\mathcal{P}}}
\newunicodechar{∣}{\ensuremath{\mid}}
\newunicodechar{∗}{\ensuremath{\ast}}
\newunicodechar{⋯}{\ensuremath{\cdots}}
\newunicodechar{≪}{\ensuremath{\ll}}
\newunicodechar{∼}{\ensuremath{\sim}}
\ifLuaTeX
  \usepackage{selnolig}  % disable illegal ligatures
\fi
\IfFileExists{bookmark.sty}{\usepackage{bookmark}}{\usepackage{hyperref}}
\IfFileExists{xurl.sty}{\usepackage{xurl}}{} % add URL line breaks if available
\urlstyle{same}
\hypersetup{
  pdftitle={The witness game beyond Erds Problem 848: higher powers, ab+c, and the third density},
  pdfauthor={Luke Culpitt (AI assisted)},
  colorlinks=true,
  linkcolor={blue},
  filecolor={Maroon},
  citecolor={Blue},
  urlcolor={Blue},
  pdfcreator={LaTeX via pandoc}}

\title{The witness game beyond Erd\H{o}s Problem 848: higher powers,
\(ab+c\), and the third density}
\author{Luke Culpitt (AI assisted)}
\date{November 2026}

\begin{document}
\maketitle

\subsection{Abstract}

Erdős Problem 848 asks for the largest A ⊆ \{1, \ldots, N\} with ab + 1
never squarefree. The answer is the class 7 mod 25. In a companion paper
{[}M{]} we proved this for every N by a finite ``witness game'' on
residue types. That paper also determined the second-best density, 1/50,
attained by a two-class set F. Here we test how far the method reaches.

\begin{itemize}
\tightlist
\item
  \textbf{Higher powers.} For ab + 1 never k-free, the maximum is a
  class mod 5\textsuperscript{k} for every k ≥ 3 and all large N, with
  an explicit threshold. For 3 ≤ k ≤ 100 this holds for every N, with
  exactly two small exceptions, at k = 3 and k = 4. For cubes the
  maximisers are listed for every N.
\item
  \textbf{Other forms.} For ab + c never squarefree, the maximum for
  large N is a root class of x² + c ≡ 0 (mod p²). This holds for 49
  squarefree c ≤ 100, and for 33 of them the second-best density is also
  determined. The counting lemmas needed for general c are proved here.
\item
  \textbf{The third density for 848.} The densest valid set not inside a
  class, F or F's mirror has density between 1/25 − 1/(4π²) = 0.014670
  and 0.014910. The lower bound is one class plus one outsider, and we
  conjecture that it is sharp.
\end{itemize}

The computations use exact clique search and CP-SAT with rounded-up
objective coefficients. All code is public.

\subsection{1. Introduction}

Call A ⊆ {[}N{]} \emph{valid} if ab + 1 is not squarefree for all a, b ∈
A, including a = b. The class C7 = \{n ≡ 7 (mod 25)\} is valid, since 7²
+ 1 = 50. In {[}M{]} we proved that it is the largest valid set for
every N, together with stability and rigidity results and the
second-best density.

The method of {[}M{]} has three ingredients:

\begin{itemize}
\tightlist
\item
  a \emph{self-forcing class}: a residue class in which ab + 1 is always
  divisible by a fixed square;
\item
  a \emph{positive local gap}: a finite game on residue types showing
  that any valid set not inside such a class is smaller by a constant
  factor;
\item
  \emph{control of large primes}: counting lemmas showing that squares
  p² with p large contribute o(N), uniformly in the witness.
\end{itemize}

This paper asks how special these are to x² + 1 and to squares. Section
2 proves the counting lemmas for ab + c with a general squarefree c ≥ 1.
Section 3 uses them, with c = 1, to bracket the third density of 848.
Section 4 treats higher powers, and Section 5 the forms ab + c.

We use the notation of {[}M{]} throughout. A \emph{type} is a set of
residue classes modulo a fixed modulus L. Type s \emph{forces} type t if
some fixed square divides ab + c for every a ∈ t and b ∈ s. A type is
\emph{self-forcing} if it forces itself. For a set A, P is the set of
types it meets, its \emph{core} K is the set of types in P forced by
every type in P, and F(K) is the set of types forcing every member of K.
Forcing is symmetric, so P ⊆ F(K).

\subsection{2. Counting lemmas for ab + c}

Fix a squarefree c ≥ 1 and a modulus L. The \emph{layer primes} are the
primes p whose squares are built into the types, through L = 36·∏ p² or
similar. Every other prime q is a \emph{luck prime}. For a luck prime q,
the event q² \textbar{} ab + c is not decided by the types and is
counted directly.

\textbf{Lemma 2.1 (large luck primes).} Let 1 ≤ b ≤ N and Y ≥ c.~Then

\#\{a ≤ N : q² \textbar{} ab + c for some prime q \textgreater{} Y\} ≤
c·2\textsuperscript{ω(b)+ω(c)+1}·(2(N + c)/Y + (N² + c)/Y²).

\emph{Proof.} Let q \textgreater{} Y have q² \textbar{} ab + c.~Then q ∤
c, since q \textgreater{} c, and so q ∤ b. Write ab + c = mq², so that 1
≤ m ≤ M = (bN + c)/Y² and mq² ≡ c (mod b).

Count the units s mod b with ms² ≡ c (mod b) prime power by prime power.
Let ℓ\textsuperscript{e} ∥ b.

\begin{itemize}
\tightlist
\item
  If ℓ ∤ c, then ℓ ∤ m, and s² ≡ cm⁻¹ has at most 2 solutions mod
  ℓ\textsuperscript{e} (at most 4 if ℓ = 2).
\item
  If ℓ \textbar{} c, then ℓ ∥ c because c is squarefree, and ℓ
  \textbar{} m. For e = 1 every unit works, which gives fewer than ℓ
  solutions. For e ≥ 2, write m = ℓm′ and c = ℓc′. Then m′s² ≡ c′ (mod
  ℓ\textsuperscript{e−1}) with c′ a unit, which has at most 4 solutions
  mod ℓ\textsuperscript{e−1}, so at most 4ℓ mod ℓ\textsuperscript{e}.
\end{itemize}

So q lies in at most c·2\textsuperscript{ω(b)+ω(c)+1} classes mod b.
Since q ≤ √((bN + c)/m), each m contributes at most
c·2\textsuperscript{ω(b)+ω(c)+1}(√((bN + c)/m)/b + 1). Summing over m ≤
M with Σ m\textsuperscript{−1/2} ≤ 2√M gives
c·2\textsuperscript{ω(b)+ω(c)+1}(2(bN + c)/(bY) + M). This is at most
the bound stated, since 1 ≤ b ≤ N. ∎

For c = 1 this is Lemma 8.3 of {[}M{]}. With Y = N\textsuperscript{0.6}
it is N\textsuperscript{0.8+o(1)}, uniformly in b ≤ N, because
2\textsuperscript{ω(b)} = N\textsuperscript{o(1)}.

\textbf{Lemma 2.2 (one witness).} Let t be a type of density
w\textsubscript{t} and let b ≤ N. For each luck prime q, let
d\textsubscript{q}(t, b) be the relative density in t of the a with q²
\textbar{} ab + c.~This density is determined by the classes of b mod L
and mod q². Then, uniformly in b ≤ N,

\#\{a ∈ t ∩ {[}N{]} : q² \textbar{} ab + c for some luck prime q\} ≤
w\textsubscript{t}·(1 − ∏\textsubscript{q} (1 − d\textsubscript{q}(t,
b)))·N + o(N),

and also ≤ w\textsubscript{t}·Σ\textsubscript{q} d\textsubscript{q}(t,
b)·N + o(N).

The densities d\textsubscript{q}(t, b) are as follows.

\begin{itemize}
\tightlist
\item
  For a luck prime q ∤ cL:

  \begin{itemize}
  \tightlist
  \item
    if q ∤ b, a must lie in the single class −cb⁻¹ mod q², so
    d\textsubscript{q} = 1/q²;
  \item
    if q \textbar{} b, then ab + c ≡ c ≢ 0 (mod q), so
    d\textsubscript{q} = 0.
  \end{itemize}
\item
  For a luck prime q \textbar{} c, the type records whether q \textbar{}
  a. Then q² \textbar{} ab + c needs q to divide exactly one of a and b.

  \begin{itemize}
  \tightlist
  \item
    If q \textbar{} b and q ∤ a, then a lies in one nonzero class mod q,
    so d\textsubscript{q} = 1/(q − 1).
  \item
    If q \textbar{} a and q ∤ b, then a/q lies in one class mod q, so
    d\textsubscript{q} = 1/q.
  \item
    In the other two cases, d\textsubscript{q} = 0.
  \end{itemize}
\item
  For a luck prime q whose root classes are recorded in the types (as in
  Sections 3 and 5):

  \begin{itemize}
  \tightlist
  \item
    if a lies in a root class R and b does not, then d\textsubscript{q}
    = 0, because ab ≡ −c (mod q²) would put b in R;
  \item
    if a lies in no root class, then −cb⁻¹ is one of the q² − 2
    remaining classes, so d\textsubscript{q} ≤ 1/(q² − 2).
  \end{itemize}
\end{itemize}

\emph{Proof.} Fix Q. For the luck primes q ≤ Q, the events are
conditions on a modulo L·∏\textsubscript{q≤Q} q². By the Chinese
remainder theorem, the count is exact up to O(1) per class, where the
implied constant depends on Q. This gives the product for q ≤ Q, and the
union bound gives the sum.

For Q \textless{} q ≤ Y with Y = N\textsuperscript{0.6}, each q
contributes at most N/q² + 1, or N/q + 1 when q \textbar{} c; but q
\textbar{} c means q ≤ c ≤ Q once Q ≥ c.~So these primes contribute at
most N/Q + π(Y). For q \textgreater{} Y use Lemma 2.1. Letting Q → ∞
slowly with N, the total error is o(N), uniformly in b. ∎

\textbf{Lemma 2.3 (large primes in x² + c).} For every Q ≥ 1,

\#\{x ≤ N : p² \textbar{} x² + c for some prime p \textgreater{} Q\} ≤
(2/Q)·N + 2π(√(N² + c)) = (2/Q)·N + O(N/log N).

\emph{Proof.} If p \textbar{} c then p² ∤ x² + c: p \textbar{} x would
give p² \textbar{} x², and c is squarefree. So p ∤ c, and x² ≡ −c (mod
p²) has at most two roots. Each root class meets {[}N{]} in at most N/p²
+ 1 elements, and p ≤ √(N² + c). Sum over p \textgreater{} Q, using
Σ\textsubscript{p\textgreater Q} 1/p² ≤ 1/Q. ∎

For c = 1, {[}M{]} used Pell's equation to get a power saving here. For
an asymptotic statement this cruder bound suffices, and it holds for
every c.

\textbf{Lemma 2.4 (joint estimate with the own witness).} Let t be a
type that is not self-forcing, and let b ∈ t. Write V\textsubscript{c} =
\{x : x² + c not squarefree\}, and ρ\textsubscript{p} for the number of
roots of x² ≡ −c (mod p²). Suppose t records, for each luck prime q
\textbar{} c, whether q \textbar{} a. Then, uniformly in b ≤ N,

\#\{a ∈ t ∩ {[}N{]} : a ∈ V\textsubscript{c}, ab + c not squarefree\} ≤
w\textsubscript{t}·(α\textsubscript{t}·β⁰\textsubscript{t} +
κ\textsubscript{t})·N + o(N),

where

\begin{itemize}
\tightlist
\item
  α\textsubscript{t} is the relative density in t of the a that have p²
  \textbar{} a² + c for some luck prime p, bounded by the union bound Σ
  ρ\textsubscript{p}/p² (a luck prime whose root classes are recorded
  contributes nothing here, since t is not self-forcing and so lies
  outside them);
\item
  β⁰\textsubscript{t} = Σ d\textsubscript{q}(t, b) over luck primes q ∤
  c;
\item
  κ\textsubscript{t} = Σ 1/p² over the luck primes p with
  ρ\textsubscript{p} = 2.
\end{itemize}

\emph{Proof.} Since t is not self-forcing and b ∈ t, the layer primes
force nothing between a and b. So a ∈ V\textsubscript{c} needs p²
\textbar{} a² + c for some luck prime p, and ab + c not squarefree needs
q² \textbar{} ab + c for some luck prime q.

A luck prime q \textbar{} c contributes nothing. The elements a and b
lie in the same type, so q divides both or neither. If q divides both,
ab + c ≡ c (mod q²) and q ∥ c.~If q divides neither, ab + c ≡ ab ≢ 0
(mod q).

Fix Q and consider p, q ≤ Q. For p ≠ q, the Chinese remainder theorem
gives the product of the densities. For p = q, a must lie in a root
class mod p² and in the class −cb⁻¹ mod p². The two root classes are
distinct, so this leaves at most one class, and contributes 1/p². Up to
the CRT errors, this gives w\textsubscript{t}(α\textsubscript{t}
β⁰\textsubscript{t} + κ\textsubscript{t})N.

For p \textgreater{} Q, drop the condition on q and use Lemma 2.3. For q
\textgreater{} Q, drop the condition on p and use the tail of Lemma 2.2.
Let Q → ∞. ∎

\textbf{Lemma 2.5 (core reduction).} Let A ⊆ {[}N{]} be c-valid, meaning
ab + c is not squarefree for all a, b ∈ A. Let P be the set of types A
meets, with core K. Suppose that for each type t there is a constant
b\textsubscript{t} with these properties:

\begin{itemize}
\tightlist
\item
  if t is self-forcing, b\textsubscript{t} is at least the luck density
  of Lemma 2.2 for every witness that does not force t;
\item
  otherwise, b\textsubscript{t} is at least the constant of Lemma 2.4.
\end{itemize}

Then

\textbar A\textbar{} ≤ U(K)·N + o(N), where U(K) = w(K) +
Σ\textsubscript{t~∈~F(K)~∖~K} b\textsubscript{t}·w\textsubscript{t}.

\emph{Proof.} As in {[}M, Lemma 8.5{]}. Every element of A is compatible
with every other. A self-forcing type t ∈ P ∖ K is not forced by some s
∈ P, so any element of A ∩ s is a witness for Lemma 2.2. A type that is
not self-forcing uses an element of itself as the witness for Lemma 2.4.
Core types count in full. Since P ⊆ F(K), summing over t ∈ P gives the
bound. The number of types is fixed, so the o(N) terms add up to o(N). ∎

Maximising U(K) over the cores that a given family of sets can have is a
finite problem. We solve it with CP-SAT {[}OR-tools{]}, with every
objective coefficient rounded up to a multiple of 10⁻⁹, so the reported
optimum is an upper bound.

The variables are as follows. A \emph{slab} is a residue class for the
layers other than 4 and 9, together with the root-class memberships
described below. A \emph{cell} is a residue mod 36, together with the
divisibility flags. Then K is given by a closed set of cells in each
slab, and F(K) by its intersections, exactly as in {[}M, Section 8.4{]}.
The code is \texttt{coregame.py}.

\subsection{3. The third density of Problem 848}

Let F = \{7 mod 25, 1 mod 4\} ∪ \{18 mod 25, 3 mod 4\} and let F′ be its
mirror image. In {[}M{]} we showed that every valid set not inside C7 or
C18 has at most N/50 + O(N\textsuperscript{0.81}) elements, and that F
and F′ attain this. We also showed that for N ≥
10\textsuperscript{23.6}, a valid set not inside a class with at least
(1/50 − 0.0008)N elements lies in F or F′. Let λ₃(N) be the largest size
of a valid A ⊆ {[}N{]} contained in none of C7, C18, F and F′.

\textbf{Theorem 3.1.} 1. lim inf λ₃(N)/N ≥ 1/25 − 1/(4π²) =
0.0146697\ldots. 2. lim sup λ₃(N)/N ≤ 0.014910.

So for all large N, a valid set with more than 0.01491N elements lies
inside C7, C18, F or F′. This sharpens the density 1/50 − 0.0008 =
0.0192 of {[}M, Theorem 8.9{]} to 0.01491, though without an explicit
threshold for N.

\emph{Proof of (1).} Take a prime x \textgreater{} 3 with x ≡ 18 (mod
25) and x ≡ 3 (mod 4), and let

A = \{x\} ∪ \{a ∈ C7 ∩ {[}N{]} : ax + 1 not squarefree\}.

Pairs inside C7 are forced by 25. Also x² + 1 ≡ 0 (mod 25), and pairs
with x hold by definition, so A is valid. It contains x ∈ C18 and
elements of C7 that are ≡ 1 and ≡ 3 (mod 4), so it lies in none of the
four sets.

For a ∈ C7 we have ax + 1 ≡ 2 (mod 25), so 25 never divides ax + 1.
Among the a ∈ C7:

\begin{itemize}
\tightlist
\item
  4 divides ax + 1 for one class of a mod 4;
\item
  9 divides it for one class mod 9;
\item
  p² divides it for one class mod p², when p ≥ 7 and p ≠ x;
\item
  x² never divides it.
\end{itemize}

So the a ∈ C7 with ax + 1 squarefree have density

(1/25)·(3/4)·(8/9)·∏\textsubscript{p≥7,~p≠x}(1 − p⁻²) = (1/(4π²))·(1 −
x⁻²)⁻¹,

using ∏\textsubscript{p≥7}(1 − p⁻²) = (6/π²)·(4/3)·(9/8)·(25/24) =
75/(8π²). Hence \textbar A\textbar/N → 1/25 − (1/(4π²))(1 − x⁻²)⁻¹ as N
→ ∞, and letting x → ∞ gives (1). The script \texttt{lam3\_lower.py}
checks the count: 0.014671 at N = 2·10⁷ with x = 100043. ∎

This is the ``one class plus one outsider'' configuration. The outsider
forces exactly a third of C7, through 4 and 9 {[}M, Lemma 2.2{]}. The
other two thirds survive only where ax + 1 picks up a square factor by
chance. Its density is the sharp constant g₁* of {[}M, Theorem C(1){]}.

\emph{Proof of (2).} Apply Lemma 2.5 with c = 1 and the following types:

\begin{itemize}
\tightlist
\item
  the class mod 25;
\item
  the cell mod 36;
\item
  membership in the classes S\textsubscript{p}\^{}± = \{a ≡
  ±ι\textsubscript{p} (mod p²)\} for p = 13 and 17, where
  ι\textsubscript{p}² ≡ −1.
\end{itemize}

Two elements of the same S\textsubscript{p}\^{}σ are forced. An element
of S\textsubscript{p}\^{}σ and one outside it never have p² \textbar{}
ab + 1, since ab ≡ −1 (mod p²) would put the second element in
S\textsubscript{p}\^{}σ too. Inside the complement of
S\textsubscript{p}\^{}±, the prime p acts as a luck prime with density
1/(p² − 2). So the luck primes are q ≥ 7, and the types outside C7 ∪ C18
∪ S₁₃\^{}± ∪ S₁₇\^{}± are in V only through primes p ≥ 29.

The constants are as follows.

\begin{itemize}
\tightlist
\item
  For self-forcing types, b\textsubscript{t} = 1 − ∏(1 − q⁻²), with the
  product over 7 ≤ q ≤ 10⁷ and the factor for p = 13, 17 replaced by 1 −
  1/(p² − 2), plus Σ\textsubscript{q\textgreater10⁷} q⁻². This is
  0.0502.
\item
  For the remaining types, b\textsubscript{t} is the Lemma 2.4 constant,
  0.0050. It was 0.0158 for the unrefined R-types of {[}M{]}.
\end{itemize}

The admissible cores are those whose F(K) has a type outside each of C7,
C18, F and F′. Maximising U over them is a problem on 225 slabs × 21
cells. \texttt{third\_cp.py} solves it to optimality and finds
0.0149093. ∎

The optimal core is the class-plus-outsider core: twelve cells of one
class, all forced by a single cell. So the game identifies the right
configuration, and the gap of 0.00024 to the construction has two
sources.

\begin{itemize}
\tightlist
\item
  Most of it comes from elements of the outsider's cell, in the other
  classes mod 25, that lie in V through primes p ≥ 29. This part shrinks
  as more classes S\textsubscript{p} are split off: the bound is
  0.0149645 with S₁₃ alone.
\item
  The rest, β/900 ≈ 0.000056, is the outsider cell itself. The game
  charges it at rate β, although the construction has a single element
  there.
\end{itemize}

Closing the second part needs a bound for pairs of sets X, Y with xy + 1
never squarefree that is better than one witness on each side.

The unrefined enumeration of {[}M, Section 8.4{]} gives only 0.0186 here
(\texttt{third.py}). In that version the R-types carry the full constant
κ = Σ\textsubscript{p≡1(4),~p≥13} p⁻² = 0.0138 from the coincidence p =
q, which the refinement removes.

\textbf{Conjecture 3.2.} λ₃(N)/N → 1/25 − 1/(4π²).

\subsection{4. Higher powers}

\subsubsection{4.1 Cubes}

Call A ⊆ {[}N{]} \emph{cube-valid} if ab + 1 is not cube-free for all a,
b ∈ A, including a = b. Write V₃ = \{x : x² + 1 not cube-free\}. Since
57² + 1 = 2·5³·13, the classes 57 and 68 mod 125 are cube-valid, playing
the role of 7 and 18 mod 25. Write c57(N) = \#\{n ≤ N : n ≡ 57 (mod
125)\}.

\textbf{Theorem 4.1.} For every N, every cube-valid A ⊆ {[}N{]} has
\textbar A\textbar{} ≤ c57(N). Equality holds exactly in the following
cases:

\begin{itemize}
\tightlist
\item
  A = C57 ∩ {[}N{]};
\item
  A = C68 ∩ {[}N{]}, when c68(N) = c57(N);
\item
  for 193 ≤ N ≤ 306, also A = \{182, 193\};
\item
  for 239 ≤ N ≤ 306, also A = \{57, 239\} and A = \{193, 239\}.
\end{itemize}

The exceptional sets work through 8 and 27: 182·193 + 1 = 3³·1301,
57·239 + 1 = 2³·13·131 and 193·239 + 1 = 2⁴·3·31². Locally the picture
is the same as for 848. The densest congruence-defined cube-valid set
not inside one class is \{57 mod 125, 1 mod 8\} ∪ \{68 mod 125, 7 mod
8\}, of density 1/500 (\texttt{wgame4.py}).

The proof is much shorter than for 848, for two reasons. The local gap
is far larger, so a single witness suffices and the full game is not
needed. And a cube q³ \textgreater{} N² + 1 can never divide ab + 1, so
no lattice lemma is needed either.

\textbf{Lemma 4.2 (one witness).} Let A be cube-valid, and suppose some
b ∈ A lies outside C57. Then

\textbar A ∩ C57\textbar{} ≤ (N/125)·(1/8 + 1/27 + β₃) + 2 + π((N² +
1)\textsuperscript{1/3}), where β₃ = Σ\textsubscript{q≥7} q⁻³ ≤
0.004726,

and the same holds for C68.

\emph{Proof.} Since b ∉ C57, 125 never divides ab + 1 for a ∈ C57. So
some prime q ≠ 5 has q³ \textbar{} ab + 1, which puts a in one class mod
q³ (if q ∤ b). For q = 2 and q = 3 this is at most a fraction 1/8 or
1/27 of C57, at a cost of 1 each. For q ≥ 7, the class mod 125q³
contains at most N/(125q³) + 1 elements up to N. Also q³ ≤ ab + 1 ≤ N² +
1. ∎

\emph{Proof of Theorem 4.1.} Let A be cube-valid and inside neither C57
nor C68. Then A has elements outside each class, so Lemma 4.2 applies to
both. The remaining elements lie in V₃ outside C57 ∪ C68, so each has p³
\textbar{} x² + 1 for some p ≡ 1 (mod 4) with p ≥ 13. There are at most
2Nσ₃ + 2π((N² + 1)\textsuperscript{1/3}) of them, where σ₃ = Σ p⁻³ over
primes p ≡ 1 (mod 4) with p ≥ 13, so σ₃ ≤ 0.000756.

\begin{itemize}
\tightlist
\item
  \textbf{Large N.} The total main term is 0.004178N against c57(N)
  \textgreater{} (N − 57)/125 ≈ 0.008N. Checking cells of ratio 1.01
  with exact prime counts to 10⁸ and Rosser--Schoenfeld above
  (\texttt{cube\_analytic.py}) gives \textbar A\textbar{} \textless{}
  c57(N) for 1.94·10⁶ ≤ N ≤ 10⁴⁰. Beyond 10⁴⁰ the error term 4π((N² +
  1)\textsuperscript{1/3}) + 4 is far below (0.008 − 0.004178)N.
\item
  \textbf{700 ≤ N ≤ 2·10⁶.} Here \texttt{cubeM.c} computes the bound
  exactly. It finds the maximum over witnesses b ∈ V₃ ∖ C57 of \#\{a ∈
  C57 ∩ {[}N{]} : ab + 1 not cube-free\}, likewise for C68, and adds
  \textbar V₃ ∩ {[}N{]}\textbar{} − c57(N) − c68(N). All three are
  nondecreasing in N. On 20 cells of ratio 1.5, the bound at the top of
  each cell stays below c57 at the bottom of the cell. At its closest it
  reaches 89\% of c57 (17 against 19 on {[}2362, 3543{]}). A brute-force
  recomputation by factorisation agrees at N = 3000.
\item
  \textbf{N ≤ 700.} \texttt{cube\_maxsets.py} lists all maximum
  cube-valid sets for every N by an exact clique search. An independent
  clique computation confirms the maximum c57(N) up to N = 40,000. ∎
\end{itemize}

For squares the large-prime events were the main difficulty. For cubes
all three ingredients are easy.

\subsubsection{4.2 General k}

Call A ⊆ {[}N{]} \emph{k-valid} if ab + 1 is never k-free, for all a, b
∈ A including a = b. Let q = 5\textsuperscript{k}, let
r\textsubscript{k} be the smaller root of x² ≡ −1 (mod q), and write
c\textsubscript{k}(N) = \#\{n ≤ N : n ≡ r\textsubscript{k} (mod q)\}.
For example, r₃ = 57, r₄ = 182, r₅ = r₆ = 1068 and r₇ = 32318.

\textbf{Theorem 4.3.} 1. For every k ≥ 3 there is an explicit
N\textsubscript{k} such that for N ≥ N\textsubscript{k}, every k-valid A
⊆ {[}N{]} not inside one of the classes ±r\textsubscript{k} (mod q) has
fewer than c\textsubscript{k}(N) elements. So for N ≥ N\textsubscript{k}
the maximum is c\textsubscript{k}(N), and only the full classes attain
it. Here N\textsubscript{k}/5\textsuperscript{k} is about 15,400 for k =
3, 555 for k = 4, at most 214 for k ≥ 5, and at most 60 for k ≥ 40. 2.
For 3 ≤ k ≤ 100 the maximum is c\textsubscript{k}(N) for every N, and
only the full classes attain it, with exactly two exceptions. For k = 3
and 193 ≤ N ≤ 306, the mixed sets of Theorem 4.1 also attain it. For k =
4 and 239 ≤ N ≤ 806, the set \{239\} also attains the maximum 1, because
239² + 1 = 2·13⁴.

\emph{Proof of part 1.} Lemma 4.2 holds with 2\textsuperscript{k},
3\textsuperscript{k} and p\textsuperscript{k} in place of 8, 27 and p³.
A witness outside C\textsubscript{rk} forces at most a fraction
Z\textsubscript{k} = 2\textsuperscript{−k} + 3\textsuperscript{−k} +
Σ\textsubscript{p≥7} p\textsuperscript{−k} of that class, plus one
element for each prime p with p\textsuperscript{k} ≤ N² + 1. The rest of
V\textsubscript{k} has at most 2Nσ\textsubscript{k} + 2π((N² +
1)\textsuperscript{1/k}) elements, where σ\textsubscript{k} = Σ
p\textsuperscript{−k} over primes p ≡ 1 (mod 4) with p ≥ 13. The main
term 2Z\textsubscript{k} + 2qσ\textsubscript{k} is largest at k = 3,
where it is 0.52, and it falls to 0.21 at k = 4 and 0.09 at k = 5.

For 3 ≤ k ≤ 40, \texttt{kfree\_analytic.py} finds N\textsubscript{k} by
checking cells of ratio 1.01. For k \textgreater{} 40 and N = λq with λ
≥ 60, the requirement reduces to λ(1 − 10⁻¹¹) − 1 \textgreater{}
4π(25.01·λ\textsuperscript{0.05}) + 4. This holds because the right side
is at most 52 for λ ≤ 10⁴, and grows only like λ\textsuperscript{0.05}.
∎

\emph{Proof of part 2.} Below N\textsubscript{k} the vertex set
V\textsubscript{k} ∩ {[}N\textsubscript{k}{]} can be listed directly, as
the roots of x² ≡ −1 (mod p\textsuperscript{k}) for the few primes with
p\textsuperscript{k} ≤ N\textsubscript{k}² + 1. Testing an edge only
needs those same primes. The vertex set is small: 1,143 vertices for k =
4, 432 for k = 5, and between 90 and 155 for 9 ≤ k ≤ 100. Above k = 30
there are no vertices outside the two classes at all.

For 5 ≤ k ≤ 100, \texttt{kfree\_exact.py} enumerates all maximum cliques
at every vertex value up to N\textsubscript{k}; the graph is constant
between consecutive vertices. For k = 4 it does this up to 30,000 and
then checks the one-witness bound on 14 cells up to N₄, where the bound
is at most a third of c₄. The case k = 3 is Theorem 4.1.

The graphs are not trivial: for k = 5 there are 844 edges between the
two classes and 5 extra vertices. An independent brute force,
factorising x² + 1 and ab + 1 directly, reproduces every maximum set for
k = 4 up to 3000 and k = 5 up to 20,000. ∎

So the extremal answer is the same for every k ≥ 3: the class
r\textsubscript{k} mod 5\textsuperscript{k}, with occasional small
exceptions at the start. Part 2 cannot be checked for all k at once.
Whether a stray vertex such as 239 falls below the second element of the
class depends on the particular roots.

The exceptions are not confined to c\textsubscript{k}(N) ≤ 1: for k = 3
they occur where c₃(N) = 2. But they do not recur. Let
μ\textsubscript{k}(N) be the size of the largest k-valid set in {[}N{]}
not inside one class. For every 5 ≤ k ≤ 300 and every N ≤
N\textsubscript{k}, μ\textsubscript{k}(N) ≤ c\textsubscript{k}(N) − 3.
The closest approach is k = 5 and N = 13,568, where a mixed pair faces a
class of 5 (\texttt{margin.py}). For 12 ≤ k ≤ 300 there is no mixed set
at all below N\textsubscript{k}: V\textsubscript{k} ∩
{[}N\textsubscript{k}{]} consists of the two classes, with no edge
between them.

\subsection{5. The forms ab + c}

Fix a squarefree c ≥ 1, and call A ⊆ {[}N{]} \emph{c-valid} if ab + c is
not squarefree for all a, b ∈ A, including a = b. A \emph{root class} is
a class x mod p² with x² + c ≡ 0 (mod p²). It is self-forcing, because
aa′ + c ≡ x² + c ≡ 0 for a, a′ in the class.

For p = 2, root classes exist exactly when c ≡ 3 (mod 4): they are 1 and
3 mod 4. For odd p ∤ c, they exist exactly when −c is a square mod p,
and then there are two classes mod p². For p \textbar{} c there are
none. Let p\textsubscript{c} be the least prime with root classes; for c
= 1 it is 5. The root classes mod p\textsubscript{c}² are the natural
candidates for the maximum. They have density 1/4 if p\textsubscript{c}
= 2 and 1/p\textsubscript{c}² otherwise.

\textbf{The game for c.} The types are built as follows.

\begin{itemize}
\tightlist
\item
  The layer primes are 2, 3 and P, where P = p\textsubscript{c} if
  p\textsubscript{c} ≥ 7 and P = 5 otherwise. A type records a mod 36
  and a mod P².
\item
  Every other prime is a luck prime. For each luck prime q \textbar{} c,
  the type records whether q \textbar{} a, as Lemmas 2.2 and 2.4
  require.
\item
  The type records membership in the root classes of the least luck
  prime p with ρ\textsubscript{p} = 2, as in Section 3.
\end{itemize}

The constants b\textsubscript{t} are those of Lemmas 2.2 and 2.4: union
bounds over all luck primes, with the q \textbar{} c terms at their
worst case for self-forcing types. Lemma 2.5 then reduces the problem to
a finite maximisation, which \texttt{game\_gen.py} solves with CP-SAT in
two steps.

\begin{enumerate}
\def\labelenumi{\arabic{enumi}.}
\tightlist
\item
  It finds the maximum-weight cliques among the types. These are the
  largest congruence-defined c-valid sets.
\item
  It maximises U(K) over the cores whose F(K) is not inside any maximum
  clique. Call this maximum V\textsubscript{c}.
\end{enumerate}

\textbf{Theorem 5.1 (computer-assisted).} Let c ≤ 100 be squarefree with
p\textsubscript{c} ≤ 7 and c ∉ \{10, 13\}; there are 49 such c.~Then:

\begin{enumerate}
\def\labelenumi{\arabic{enumi}.}
\tightlist
\item
  The maximum cliques of the game are exactly the two root classes mod
  p\textsubscript{c}² (mod 4 if p\textsubscript{c} = 2).
\item
  V\textsubscript{c} \textless{} 1/p\textsubscript{c}²
  (V\textsubscript{c} \textless{} 1/4 if p\textsubscript{c} = 2).
\item
  For all sufficiently large N, the largest c-valid sets in {[}N{]} are
  the root classes mod p\textsubscript{c}² of largest size, and only
  these.
\end{enumerate}

\emph{Proof.} Parts 1 and 2 are the computation. For part 3, let A be
c-valid and not inside a root class mod p\textsubscript{c}². Then the
type set of A is not inside a maximum clique, so neither is F(K). By
Lemma 2.5, \textbar A\textbar{} ≤ (V\textsubscript{c} + o(1))N. Since
V\textsubscript{c} is strictly below the density of the root classes,
this is smaller than the root classes once N is large. A c-valid set
inside a root class has at most its size, with equality only for the
full class. ∎

For 33 of these c, the best congruence-defined c-valid set not inside a
root class has density exactly V\textsubscript{c}. So V\textsubscript{c}
is the second-best density.

\begin{longtable}[]{@{}lll@{}}
\toprule\noalign{}
max density & second-best density & c \\
\midrule\noalign{}
\endhead
\bottomrule\noalign{}
\endlastfoot
1/4 (classes mod 4) & 1/9 & 3, 11, 23, 47, 51, 59, 71, 83, 87, 95 \\
1/4 & 1/18 & 43, 67, 79 \\
1/9 (classes mod 9) & 1/12 & 2, 14, 26, 38, 62, 74, 86 \\
1/9 & 1/18 & 5, 17, 29, 41, 53, 65, 89 \\
1/25 (classes mod 25) & 1/50 & 1 \\
1/25 & 3/100 & 6, 34, 46 \\
1/49 (classes mod 49) & 3/196 & 82 \\
1/49 & 1/98 & 97 \\
\end{longtable}

For the other 16 values of c (7, 15, 19, 21, 31, 33, 35, 39, 55, 61, 66,
69, 73, 77, 91, 94), the best congruence set and V\textsubscript{c}
still differ, so the second-best density is only bracketed
(\texttt{gen\_k1.jsonl}).

The refinements matter. Without the root-class split and the
divisibility flags, the game fails for c = 34, 61, 66, 69 and 94: for c
= 34 it gives V = 0.0427 against 1/25. Each of these has a luck prime
dividing c, such as 17 \textbar{} 34, or many small primes p with
ρ\textsubscript{p} = 2. With the refinements all five pass. With c = 1
the computation reproduces V = 1/50 of {[}M{]}.

Exact clique computations (\texttt{brute\_c.py}) agree with the root
class for every N up to 400--800 for c = 2, 3, 6, 10, 13, 21 and 34.

Two groups of c ≤ 100 are left open.

\begin{itemize}
\tightlist
\item
  \textbf{c = 10 and 13.} The class is mod 49, and the game with one
  split prime falls just short: V₁₃ = 0.020465 against 1/49 = 0.020408.
  A second split prime should settle them, but it exceeds our computing
  budget.
\item
  \textbf{The ten c ≤ 100 with p\textsubscript{c} ≥ 11:} 22, 30, 37, 42,
  57, 58, 70, 78, 85 and 93. The model has p\textsubscript{c}² slabs,
  and we have not run it.
\end{itemize}

Theorem 5.1 gives no explicit threshold for N, since the o(N) terms of
Section 2 are not made explicit.

\subsection{6. Open problems}

\begin{enumerate}
\def\labelenumi{\arabic{enumi}.}
\tightlist
\item
  \textbf{The third density.} Prove Conjecture 3.2. This needs a
  two-sided bound for pairs X, Y with xy + 1 never squarefree that
  improves on one witness per side.
\item
  \textbf{All k.} Does Theorem 4.3(2) hold for every k? For large k, a
  mixed set below N\textsubscript{k} needs 13\textsuperscript{k} or
  17\textsuperscript{k} to divide x² + 1 or ab + 1 with x, a, b below
  about 60·5\textsuperscript{k}. Heuristically this happens for only
  finitely many k.
\item
  \textbf{All c.} Is the root class mod p\textsubscript{c}² extremal for
  every squarefree c, including c = 10, 13 and those with
  p\textsubscript{c} ≥ 11? Is there a proof that does not need a
  computation for each c? Non-squarefree c brings new self-forcing
  classes, such as 0 mod q when q² \textbar{} c, and is not covered
  here.
\item
  \textbf{Effective thresholds} for Theorems 3.1 and 5.1.
\end{enumerate}

\subsection{Appendix: reproducibility}

All code is in the repository {[}R{]}, in the folders
\texttt{third\_density/}, \texttt{cubefree/}, \texttt{kfree/} and
\texttt{general\_c/}. The CP-SAT computations use OR-tools 9.15.

\begin{longtable}[]{@{}
  >{\raggedright\arraybackslash}p{(\columnwidth - 4\tabcolsep) * \real{0.3333}}
  >{\raggedright\arraybackslash}p{(\columnwidth - 4\tabcolsep) * \real{0.3333}}
  >{\raggedright\arraybackslash}p{(\columnwidth - 4\tabcolsep) * \real{0.3333}}@{}}
\toprule\noalign{}
\begin{minipage}[b]{\linewidth}\raggedright
Claim
\end{minipage} & \begin{minipage}[b]{\linewidth}\raggedright
Computation
\end{minipage} & \begin{minipage}[b]{\linewidth}\raggedright
Result
\end{minipage} \\
\midrule\noalign{}
\endhead
\bottomrule\noalign{}
\endlastfoot
Theorem 3.1(1) & \texttt{lam3\_lower.py}, exact count & 0.014671 at N =
2·10⁷ \\
Theorem 3.1(2) & \texttt{third\_cp.py\ 13,17\ product} & max U =
0.0149093, OPTIMAL; 0.0149645 with S₁₃ only; 0.0186 unrefined
(\texttt{third.py}) \\
Theorem 4.1 & \texttt{cube\_maxsets.py} (N ≤ 700),
\texttt{cube\_small.py} (to 40,000), \texttt{cubeM.c} (700 ≤ N ≤ 2·10⁶),
\texttt{cube\_analytic.py} (N ≥ 1.94·10⁶) & as stated \\
Theorem 4.3 & \texttt{kfree\_analytic.py}, \texttt{kfree\_exact.py},
\texttt{run\_all.py}, \texttt{run\_big.py}, \texttt{margin.py} & no
exceptions for 5 ≤ k ≤ 100; margin at least 3 for 5 ≤ k ≤ 300 \\
Theorem 5.1 & \texttt{game\_gen.py\ 1\ c} for each c; output
\texttt{gen\_k1.jsonl} & V\textsubscript{c} \textless{}
1/p\textsubscript{c}² for 49 values; V\textsubscript{1} = 1/50 \\
Small-N checks for ab + c & \texttt{brute\_c.py} & root class maximal
for c = 2, 3, 6, 10, 13, 21, 34 up to N = 400--800 \\
\end{longtable}

\subsection{References}

{[}M{]} L. Culpitt, Stability, rigidity and exact extremisers for Erdős
Problem 848, companion paper (2026).

{[}R{]} L. Culpitt, Erdos-848, https://github.com/MeaCulpitt/Erdos-848
(2026).

{[}B{]} T. F. Bloom, Erdős Problem \#848,
https://www.erdosproblems.com/848.

{[}OR-tools{]} L. Perron and F. Didier, CP-SAT, Google OR-Tools,
https://developers.google.com/optimization.

{[}RS{]} J. B. Rosser and L. Schoenfeld, Approximate formulas for some
functions of prime numbers, Illinois J. Math. 6 (1962), 64--94.

\subsection{Acknowledgements}

The arguments and code were developed with AI assistance (Claude) and
independently re-run. They have not yet been reviewed by a human expert.

\end{document}
